\section{Related Work}
\label{sec:related}

\paragraph{RLVR and its control groups.}
GRPO~\cite{grpo2024} and its descendants~\cite{deepseekr1} optimize against
an executable checker rather than a learned reward model, which removes reward
hacking of the preference-model kind and makes the training signal cheap on
program-annotated tasks. Empirical reports pair RLVR against the base policy
and a matched SFT run; stronger evaluations add rejection-sampling SFT or
iterated self-training~\cite{star2022,rft2023,restem2023}, which use the same
verifier at training time without policy-gradient updates. Every arm in these
control groups is a learning arm whose accuracy grows with the demonstration
budget. Our contribution is not a new algorithm but the observation that adding
one non-learning arm changes the sign of the reported effect.

Recent critical work on R1-Zero-style training has begun to audit the recipe
itself rather than its reported margins. Dr.\ GRPO~\cite{drgrpo2025} identifies
an optimization bias in GRPO's length normalization, and
SimpleRL-Zoo~\cite{simplerlzoo2025} shows how sharply the outcome of zero-RL
training depends on the base model across ten backbones. Both findings concern
what happens \emph{inside} the RL arm. Our question is orthogonal and
complementary: holding the RL arm fixed, what does the control group have to
contain before its reported margin means anything? A paper can adopt Dr.\ GRPO's
unbiased estimator and still report a gain that a non-learning arm erases.

\paragraph{Does RL exceed the base model's reach?}
A parallel line asks whether RLVR expands a model's capability or merely
sharpens its sampling distribution, e.g.\ by comparing pass@$k$ before and
after RL~\cite{yue2025limit} or by auditing where reported gains come
from~\cite{sober2025}. That question is about the policy; ours is about the
comparison. The two can agree in direction while differing in what would refute
them: our claim would be refuted by a strong solver that still loses to GRPO,
not by a demonstration that RL narrows the output distribution.

\paragraph{Weak baselines as a source of illusory progress.}
The stance we adopt has an explicit lineage, which the conference version of
this paper invoked without citing. In recommendation, Ferrari Dacrema et
al.~\cite{dacrema2019} reproduced only 7 of 18 recent neural methods and found
6 of those 7 beaten by properly tuned classical baselines. In deep metric
learning, Musgrave et al.~\cite{musgrave2020} found a decade of reported gains
largely attributable to experimental protocol. In neural
retrieval,~\cite{lin2018hype} made the same argument about comparisons against
untuned BM25. Deep RL has its own version of the problem: Henderson et
al.~\cite{henderson2018matters} showed that reported margins can be smaller
than the variation induced by random seeds, and Agarwal et
al.~\cite{agarwal2021precipice} showed that point estimates over few runs
routinely misrepresent aggregate performance.

We add two elements specific to RLVR. First, the baseline we restore is
\emph{non-learning}, so---unlike a tuned BM25 or a tuned matrix
factorization---it cannot be strengthened by spending more of the demonstration
budget, which is what makes it bind precisely in the scarce regime where RLVR
is argued to be valuable. Second, the decision rule is fixed by
pre-registration before the second task is run, so the comparison cannot be
retrofitted to the result.

\paragraph{Test-set reuse and selection.}
Our design reserves a frozen test set and evaluates each arm on it exactly
once, which is the operational form of a long-standing concern about adaptive
analysis. Dwork et al.~\cite{dwork2015holdout} give the theory of how repeated
querying of a holdout destroys its validity; Recht et
al.~\cite{recht2019} give the empirical counterpart on ImageNet. Our
Section~\ref{sec:results} reports three instances of the same phenomenon inside
a single study, and Section~\ref{sec:statistics} audits our own decision
quantity for the selection bias it could itself introduce. For multiple
comparisons we use Holm's sequentially rejective procedure~\cite{holm1979},
and for the interpretation of non-significant per-group results we follow
Card et al.~\cite{card2020power} on the power that small NLP test sets actually
afford.

\paragraph{Symbolic and program-based numerical reasoning.}
Program-of-Thought~\cite{pot2022} and PAL~\cite{pal2023} delegate arithmetic
to an interpreter, which is what makes verifiable rewards natural on these
tasks. The same property makes a purely deterministic solver feasible:
FinQA~\cite{finqa2021}, ConvFinQA~\cite{convfinqa2022} and
TAT-QA~\cite{tatqa2021} annotate answers with short executable programs over
semi-structured tables, and BizBench~\cite{bizbench2024} packages
CodeFinQA and CodeTAT-QA in that form.

The observation that policy gradients can be trained against an executable
checker on exactly this kind of table-program task is not new, and predates
RLVR by most of a decade. Neural Symbolic Machines~\cite{nsm2017} trained a
sequence-to-program model against a Lisp interpreter with REINFORCE under weak
supervision; MAPO~\cite{mapo2018} added a memory buffer of verified
trajectories to reduce gradient variance, on WikiTableQuestions~\cite{wtq2015}
and WikiSQL. The control-group problem we describe is visible in that
literature too, where deterministic semantic parsers remained competitive
baselines. Later, Binder~\cite{binder2023} reached state of the art on the same
table-QA benchmarks with no training at all, by binding a language model into a
symbolic program---the strongest existing evidence that on table-program tasks
the non-learning route is a serious competitor rather than a straw arm. Our
solver is a port of classical label-alignment and operator-template machinery,
not a novel method, which is exactly why its omission from RLVR control groups
is hard to justify.

\paragraph{Pre-registration and frozen evaluation.}
Pre-registration, held-out test sets touched once, and explicit deviation logs
are standard in fields where negative results must be credible. They are rare
in ML practice but disproportionately valuable for a paper whose main result is
a null: without a decision rule fixed in advance, a reader cannot distinguish
``the effect was absent'' from ``the analysis was run until the effect was
absent.'' Appendix~\ref{app:reproduce} states what was frozen when, which
comparisons are confirmatory and which are exploratory, and where the full
protocol and its logged deviations can be read.
